In fixed linear coordinates on $\mathbb R^d$ let
$\delta_r v=(r^{k_i}v_i)_{i=1}^d$, with arbitrary fixed real
$k_i>1$. A thin anisotropic open box need not contain any straight
ray, even when all of its small dilates are available. For example, let
$d=2$, $k=(2,4)$ and $U=(1,2)^2$. A displacement $(X,Y)\in\delta_r U$
then satisfies
\[
 r^2<X<2r^2,\quad r^4<Y<2r^4,\qquad
 \tfrac14 X^2<Y<2X^2 .
\]
No fixed straight ray contains such displacements arbitrarily near zero:
on a ray with $X>0$ the ratio $Y/X$ is constant, whereas here it is
positive and tends to zero. To apply a one-dimensional differentiation
theorem we therefore construct a curve whose increments can be assembled
from the available boxes.

\begin{lemma}[One-sided infinite descent]\label{one-sided}
For a finite right-continuous function $g$ on an open interval $I$,
the set of $t$ with
$\limsup_{h\downarrow0}(g(t+h)-g(t))/h=-\infty$ has outer measure zero.
\end{lemma}
\begin{proof}
Write $D^+g(t)=\limsup_{h\downarrow0}(g(t+h)-g(t))/h$. For $j,n\ge1$ let
\[
 E_{j,n}=\left\{t:
 \begin{array}{l}
 [t,t+1/j]\subset I,\quad |g(t)|\le n,\\
 g(t+h)-g(t)\le-h\quad(0<h<1/j)
 \end{array}\right\}.
\]
Right continuity permits testing $h\in\mathbb Q$ and gives
$E_{j,n}$ Borel. Also $\{D^+g=-\infty\}\subset\bigcup_{j,n}E_{j,n}$.
Cover $I$ by countably many intervals $J$ of diameter $<1/j$. For $E=J\cap E_{j,n}$,
\[
 s,t\in E,\ s>t\ \Longrightarrow\ g(s)\le g(t)-(s-t);
 \qquad G(t)=\sup(\{-n\}\cup g(E\cap[t,\infty))).
\]
Thus $G$ is bounded, nonincreasing and $G|_E=g|_E$. At almost every
$t\in E$, $t$ is a density point of $E$ and $G'(t)\in\R$. Choose
$t_i\in E$ decreasing strictly to $t$. Then
\[
 D^+g(t)\ge\lim_{i\to\infty}\frac{g(t_i)-g(t)}{t_i-t}
       =\lim_{i\to\infty}\frac{G(t_i)-G(t)}{t_i-t}=G'(t)>-\infty .
\]
Consequently $|E\cap\{D^+g=-\infty\}|^*=0$ for every $j,n,J$.
\end{proof}

\begin{lemma}[Anisotropic descent boxes]\label{box}
For continuous $V$ on an open set, the following base points $z$
form a set of Lebesgue outer measure zero. There are a nonempty
open box $U$, $c>0$ and $r_0>0$ such that for $0<r<r_0$,
\[
          V(z+\delta_r b)\le V(z)-cr \qquad(b\in U).
\]
Countably many adapted coordinate systems, translations, and
weight lists may be used. This includes arbitrarily narrow boxes.
\end{lemma}
\begin{proof}
By shrinking use rational $c,r_0$ and
$U=a+(-e,e)^d$ with rational $a,e$. Work on buffered subcharts,
extending $V$ continuously locally via cutoffs when needed.
With $\mathcal R=(0,r_0)\cap\mathbb Q$ and $\mathcal B=U\cap\mathbb Q^d$,
continuity gives the Borel test
\[
 E=\bigcap_{r\in\mathcal R,\ b\in\mathcal B}
       \{z: V(z+\delta_r b)-V(z)\le-cr\}.
\]
For $\kappa=\sum_i k_i$ define
\[
 d_k(z,w)=\max_i |z_i-w_i|^{1/k_i},\quad B_r(z)=z+\delta_r(-1,1)^d,\quad
 |B_{2r}(z)|=2^\kappa|B_r(z)|.
\]
Each coordinate satisfies the triangle inequality since $1/k_i<1$.
Lebesgue differentiation for this doubling ball basis gives, for almost every $z\in E$,
\[
 \lim_{r\downarrow0}\frac{|B_r(z)\setminus E|}{2^d r^\kappa}=0 .
\]
For fixed $\epsilon,C>0$, a center $w\in B_{Cr}(z)$ satisfies
$B_{\epsilon r}(w)\subset B_{(C+\epsilon)r}(z)$. At a density point
\begin{equation}\label{density-hit}
 |B_{(C+\epsilon)r}(z)\setminus E|=o(r^\kappa)
       < 2^d\epsilon^\kappa r^\kappa=|B_{\epsilon r}(w)|.
\end{equation}
The threshold on $r$ is independent of $w$, so the small box meets $E$.

Choose $1/\min_i k_i<\alpha<1$. In $(0,1)$ place one atom of
vector mass $\delta_{r_j}a$, $r_j=2^{-\alpha j}$, at each dyadic
rational of exact denominator $2^j$, $j\ge1$. If $N_j(t)$ counts those
rationals in $(0,t]$, the cumulative sum $D$ is given coordinatewise by
\begin{equation}\label{clock-series}
 D_i(t)=a_i\sum_{j\ge1}N_j(t)2^{-\alpha j k_i},\qquad
 |a_i|\sum_{j\ge1}2^{j-1}2^{-\alpha j k_i}<\infty .
\end{equation}
The bound is uniform in $t$, so $D$ is right-continuous. For
$0<t<t+h<1$ let $J$ be the least order of an
atom in $(t,t+h]$. There is exactly one atom at that order, and
$h\le 2^{1-J}$ since the interval contains no coarser grid point.
Put $n_m=N_{J+m}(t+h)-N_{J+m}(t)$. The spacing of the atoms gives
\[
 n_0=1,\qquad n_m\le1+h2^{J+m-1}\le 1+2^m\le 3\cdot 2^m .
\]
Fix $L$ large enough that for every $i$ the discarded tail satisfies
\[
 \begin{aligned}
 |\mathrm{tail}_i|
 &\le |a_i|\sum_{m>L}n_m r_{J+m}^{k_i}
 \le 3|a_i|r_J^{k_i}\sum_{m>L}2^{-(\alpha k_i-1)m}
 <e r_J^{k_i}/4 .
 \end{aligned}
\]
Reorder the retained scales $s_1,\ldots,s_N$ with the largest last. Then
\[
 N\le3\sum_{m=0}^L2^m,\qquad 2^{-\alpha L}r_J\le s_\ell\le s_N=r_J,\qquad
 \left|\sum_{\ell\le j}a_i s_\ell^{k_i}\right|\le N|a_i|r_J^{k_i}.
\]

At a density point $z\in E$, define $q_0=z$,
$q_j=z+\sum_{\ell=1}^j\delta_{s_\ell}a$, and choose a fixed $\epsilon>0$ with
\[
 2\epsilon^{k_i}<\tfrac14 e\,2^{-\alpha L k_i}\quad(1\le i\le d).
\]
The preceding partial-sum bound gives $q_j\in B_{Cr_J}(z)$ for a fixed $C$.
By \eqref{density-hit}, for sufficiently small $r_J$ choose
\[
 \begin{aligned}
 p_0&=q_0,\quad p_j\in E\cap B_{\epsilon r_J}(q_j)\quad(0<j<N),\\
 p_N&=z+D(t+h)-D(t),\quad \eta_j=p_j-q_j,\\
 \eta_0&=0,\quad |\eta_{j,i}|<(\epsilon r_J)^{k_i}\ (0<j<N),\quad
 \eta_N=\mathrm{tail}.
 \end{aligned}
\]
All $p_j$ for $j<N$ belong to $E$ by \eqref{density-hit}. For $j<N$,
\[
 |(p_j-p_{j-1})_i-s_j^{k_i}a_i|
 \le 2(\epsilon r_J)^{k_i}< e s_j^{k_i}.
\]
For $j=N$ the corresponding bound is
$e r_J^{k_i}/4+(\epsilon r_J)^{k_i}<e s_N^{k_i}$.
Thus each increment lies in its own admissible box. The final point need
not belong to $E$: only starting points enter the descent hypothesis.
Summing the $N$ inequalities gives
\begin{equation}\label{clock-sum}
 V(z+D(t+h)-D(t))-V(z)
 =\sum_{j=1}^N (V(p_j)-V(p_{j-1}))
 \le -c\sum_{j=1}^N s_j\le -c r_J .
\end{equation}
For each fixed $t$, $J\to\infty$ as $h\downarrow0$ and
$r_J/h\ge 2^{(1-\alpha)J-1}\to\infty$.

For every translate $z_0+D(t)$ the function $V(z_0+D(t))$ is
right-continuous and finite on a buffered chart (or use the local
extension). Wherever $z=z_0+D(t)$ is a density point of $E$ its
upper right Dini derivative is $-\infty$. Lemma~\ref{one-sided}
excludes this for almost every $t$. For the more general
Denjoy--Young--Saks alternatives see \cite[Remark 2.10]{rifford}.
Let $E^{(1)}$ denote these density points of $E$. The incidence function
$(z_0,t)\mapsto\mathbf1_{E^{(1)}}(z_0+D(t))$ is Borel. Its time integral
vanishes for each $z_0$ by the preceding derivative estimate. Tonelli's
theorem and translation invariance now give
\begin{equation}\label{incidence}
 \begin{aligned}
 0&=\int_{\R^d}\int_0^1 \mathbf1_{E^{(1)}}(z_0+D(t))\,dt\,dz_0\\
  &=\int_0^1 |E^{(1)}-D(t)|\,dt=|E^{(1)}|.
 \end{aligned}
\end{equation}
The differentiation theorem gives $|E\setminus E^{(1)}|=0$.
Taking the countable union over the rational choices finishes the proof.
\end{proof}

\begin{lemma}[A common exceptional set for fixed planes]\label{planes}
Fix a continuous $V:M\to\R$ and a countable collection of smooth charts.
In each chart prescribe countably many ordered pairs $e,f$ of distinct
coordinate axes and pairs of weights $k,l>1$. There is a Borel set $N$
of zero smooth volume containing every point $y$ at which one of these
choices gives an open rectangle $U$ and constants $c,r_0>0$ satisfying
\[
 V(y+r^k a e+r^l b f)\le V(y)-cr
 \quad((a,b)\in U,\ 0<r<r_0).
\]
Addition is in the chosen chart. The rectangle and constants may depend
on the point and on any curve witnessing the inequality.
\end{lemma}
\begin{proof}
Fix a chart $\chi$, two axes and weights $(k,l)$. Cover its target by
countably many product boxes $B$ with compact closure inside it. With a
smooth cutoff equal to one near $\overline B$, extend the coordinate
value to $\widetilde V\in C^0(\R^5)$. At a witness point shrink $U,r_0$
to stay in $B$. For fixed bounded rational $U$ and positive rational $c,r_0$,
\[
 E_\lambda=B\cap
 \bigcap_{\substack{(a,b)\in U\cap\mathbb Q^2\\r\in(0,r_0)\cap\mathbb Q}}
 \{z:\widetilde V(z+r^kae+r^lbf)-\widetilde V(z)\le-cr\}.
\]
Continuity makes this Borel test equivalent to testing all real parameters.
In the following display $V$ denotes the extended value. In coordinates $(w,z)\in\R^3\times\R^2$,
\[
 \begin{aligned}
 (E_\lambda)_w&\subset\{z:V(w,\cdot)\text{ satisfies Lemma~\ref{box} at }z\},\\
 |E_\lambda|_5&=\int_{\R^3}|(E_\lambda)_w|_2\,dw=0,\\
 N&=\bigcup_{\lambda}\chi_\lambda^{-1}(E_\lambda),\quad
 \mu(N)\le\sum_{\lambda}\int_{E_\lambda}\rho_\lambda(\zeta)\,d\zeta=0 ,
 \end{aligned}
\]
where $\rho_\lambda$ is the coordinate density. The union ranges over a
countable set; every admissible rectangle contains a rational test rectangle.
\end{proof}
